% Source-aware proof for the exact candidate statement
% p063_curve_cld_monotone_regularization_statement_20260827.tex.
% Date: 27 August 2026.
% This packet is isolated from the live book and has not received a final
% production-Rethlas verdict.

\paragraph{Imported facts and conventions.}
We use Thuillier's potential theory on smooth strictly analytic curves.
In particular, lisse functions factor through finite semistable skeleta by
\cite[Proposition~3.2.4]{Thu05}, lisse subharmonicity is characterized by
the nonnegativity of the graph Laplacian by
\cite[Proposition~3.4.4]{Thu05}, and every subharmonic function is the
pointwise infimum of a decreasing net of lisse subharmonic functions on
each relatively compact open set by \cite[Theorem~3.4.2]{Thu05}.  We use
the local semistable structure and the insertion of finitely many type
$2$ or type $3$ points into a skeleton from
\cite[\S2.2.3 and Lemma~2.2.22(ii)]{Thu05}.

At a type-$2$ point \(v\), define \(\operatorname{div}_v(h)\) to be the
divisor whose coefficient at the residual point \(p_\xi\) is the outgoing
slope \(d_\xi h(v)\).  With this definition the smooth image is the
principal subspace.  If the source uses the globally opposite divisor-sign
convention, both residual maps below are multiplied by \(-1\); the
principal subspace is invariant under this change, and the prescribed
outgoing slopes are unchanged.  The two exact residual-divisor sequences
used below are proved in \cite[Lemma~2.3.22]{Thu05}.  The identification of the
subsheaf generated by local analytic units with the sheaf of ordinary
Chambert--Loir--Ducros smooth harmonic functions is
\cite[Remark~15.10]{GJR26}.  This is the ordinary smooth sheaf, not the
weakly smooth sheaf of harmonic tropicalizations.

\begin{proof}[Proof of \cref{lem:p063-local-tree-affinoid}]
Replace \(C\) first by the smooth projective compactification of the
connected component containing \(x\).  The original analytification is an
open subspace of the compactification, so \(W\) is unchanged for the local
question.  If \(x\) is of type \(2\), \(3\), or \(4\), choose a simple
connected open neighbourhood \(U\subset W\) of \(x\) such that
\[
                         \mathbb H(U):=U\setminus U(K)
\]
is an \(\mathbf R\)-tree.  Such neighbourhoods form a fundamental system by
\cite[Corollary~4.27 and Proposition~5.9]{BPR13Structure}.

Strict affinoid neighbourhoods form a fundamental system on a smooth strict
analytic curve; see the proof of \cite[Corollary~2.1.18]{Thu05}.  Hence
there is a strict affinoid neighbourhood \(A\) of \(x\) with \(A\subset U\).
Let \(Y\) be the connected component of \(A\) containing \(x\).  A strict
affinoid algebra has finitely many connected components, each cut out by an
idempotent.  Thus \(Y\) is a connected strict affinoid domain.  Since
\(x\in A^\circ\) and \(Y\) is open and closed in \(A\), we also have
\(x\in Y^\circ\).

By \cite[Theorem~2.3.8]{Thu05}, after a finite separable extension of \(K\),
the domain \(Y\) admits a strictly semistable formal model.  The extension is
trivial because \(K\) is algebraically closed.  Its skeleton \(\Gamma_0\) is
a finite connected metric graph contained in
\(\mathbb H(Y)\subset\mathbb H(U)\).  An \(\mathbf R\)-tree contains no
embedded circuit, so \(\Gamma_0\) is a finite metric tree.

If \(x\) is of type \(1\), use
\cite[Corollary~4.27(1)]{BPR13Structure} to identify an open neighbourhood
\(U\subset W\) with the standard open unit ball.  In a coordinate \(T\),
choose \(R\in|K^\times|\) with \(|T|_x<R<1\) and put
\[
                         Y=\{|T|\leq R\}.
\]
Then \(Y\) is a connected strict affinoid neighbourhood of \(x\), and its
one-point Gauss skeleton \(\Gamma_0\) is a finite tree.

Let now \(f\) be lisse on a neighbourhood of \(Y\).  By
\cite[Proposition~3.2.4]{Thu05}, after a strictly semistable refinement there
are a retraction \(\tau:Y\to\Gamma\) and a continuous piecewise affine
function \(F\) on \(\Gamma\) such that \(f|_Y=F\circ\tau\).  Take a common
semistable refinement with \(\Gamma_0\) and insert all type-\(2\)
breakpoints; retain every type-\(3\) breakpoint as a metric subdivision
point.  The refined skeleton is still contained in
\(\mathbb H(Y)\subset\mathbb H(U)\); in the type-\(1\) case it lies in the
norm-point tree of the open ball.  Hence it is a finite metric tree.  After
the indicated subdivision, \(F\) is affine on each edge, and every added
interior vertex is of type \(2\) or type \(3\).
\end{proof}

\begin{proof}[Proof of
\cref{lem:p063-tree-rational-strictification}]
If \(\Gamma\) has no edge, take \(G=F\).  Assume henceforth that it has an
edge.  Choose \(b\in\partial\Gamma\), regard \(b\) as the root, and put
\[
                 I=V(\Gamma)\setminus\partial\Gamma.
\]
Define \(\mu:V(\Gamma)\to\mathbf Z\) by
\[
 \mu(v)=
 \begin{cases}
  1,&v\in I,\\
  0,&v\in\partial\Gamma\setminus\{b\},\\
  -\#I,&v=b.
 \end{cases}
\]
Thus \(\sum_v\mu(v)=0\).

For an edge \(e\), let \(\Gamma_e\) be the component of
\(\Gamma\setminus e^\circ\) not containing \(b\), and orient \(e\) from
\(\Gamma_e\) towards \(b\).  Prescribe on \(e\) the slope
\[
                         s_e=\sum_{v\in\Gamma_e}\mu(v)
\]
and normalize the resulting piecewise affine function \(\psi\) by
\(\psi(b)=0\).  The absence of circuits makes these prescriptions
compatible and determines \(\psi\).  If \(v\ne b\), the outgoing slope on
the edge from \(v\) to its parent is
\[
 \mu(v)+\sum_{\substack{e\text{ joining }v\\
                        \text{to a child subtree}}}s_e,
\]
whereas the outgoing slope from \(v\) towards a child subtree is
\(-s_e\).  Therefore
\[
                         \Delta_\Gamma\psi(v)=\mu(v).
\]
The same equality at \(b\) follows from \(\sum_v\mu(v)=0\).  In
particular,
\begin{equation}
\label{eq:p063-unit-interior-laplacian}
                         \Delta_\Gamma\psi(v)=1
                         \qquad(v\in I).
\end{equation}

Choose \(\delta>0\) so small that
\[
                   \delta\lVert\psi\rVert_\infty
                   <\frac{\varepsilon}{2}.
\]
Then \(H=F+\delta\psi\) satisfies
\[
                         \Delta_\Gamma H(v)\geq\delta
                         \qquad(v\in I)
\]
by \eqref{eq:p063-graph-subharmonicity} and
\eqref{eq:p063-unit-interior-laplacian}.  Denote the slope of \(H\) on an
oriented edge \(e\) by \(a_e\), and choose \(q_e\in\mathbf Q\) with
\[
                         |q_e-a_e|<\eta.
\]
There is a unique continuous function \(G\) with \(G(b)=H(b)\) and slope
\(q_e\) on every oriented edge.  If \(d\) is the largest valency of a
vertex and \(L\) is the sum of all edge lengths, then
\[
 |\Delta_\Gamma G(v)-\Delta_\Gamma H(v)|\leq d\eta,
 \qquad
 \lVert G-H\rVert_\infty\leq L\eta.
\]
The second estimate follows by integration along the unique path from
\(b\).  Taking
\[
              d\eta<\frac{\delta}{2},
              \qquad
              L\eta<\frac{\varepsilon}{2}
\]
proves \eqref{eq:p063-graph-rational-error} and
\eqref{eq:p063-graph-strict-mass}.
\end{proof}

\begin{proof}[Proof of \cref{lem:p063-residual-divisor-moving}]
After replacing a denominator-clearing integer by a sufficiently large
multiple, choose \(N\) so that \(NB\) is integral and
\[
                         \deg(NB)\geq2g+1.
\]
Put \(L=\mathcal O_{C_0}(NB)\).  For every \(p\in C_0\), Serre duality
gives
\[
 H^1(C_0,L(-p))
 \simeq
 H^0\bigl(C_0,\omega_{C_0}\otimes L^{-1}(p)\bigr)^\vee.
\]
The line bundle on the right has degree
\[
                   2g-1-\deg L<0,
\]
so \(H^1(C_0,L(-p))=0\).  The exact sequence
\[
 0\longrightarrow L(-p)\longrightarrow L
   \longrightarrow L|_p\longrightarrow0
\]
therefore makes \(H^0(C_0,L)\to L|_p\) surjective.  Its kernel is a
proper linear subspace.

For \(p\in M\), choose a nonzero \(k_0\)-linear functional
\(\lambda_p:H^0(C_0,L)\to k_0\) representing evaluation after a
trivialization of \(L|_p\).  The polynomial
\[
                         \prod_{p\in M}\lambda_p
\]
is nonzero.  A nonzero polynomial on a finite-dimensional vector space
over an infinite field cannot vanish everywhere; this follows by
induction on the dimension from the corresponding assertion in one
variable.  We may choose \(s\in H^0(C_0,L)\) with \(s(p)\ne0\) for every
\(p\in M\).  Its zero divisor \(A=(s)_0\) is effective, is disjoint from
\(M\), and satisfies
\[
                         \mathcal O_{C_0}(A)\simeq L
                         \simeq\mathcal O_{C_0}(NB).
\]
Thus \eqref{eq:p063-moving-divisor} holds.  Moreover,
\[
                         N\left(B-\frac1N A\right)=NB-A
\]
is principal.
\end{proof}

\begin{proof}[Proof of \cref{lem:p063-principal-residual-germ}]
Let \(\mathcal H_v\) be the stalk of Thuillier harmonic germs and let
\(\mathcal F_v\) be the stalk generated over \(\mathbf R\) by germs
\(\log|f|\), with \(f\) an invertible analytic function.  The proof of
\cite[Lemma~2.3.22]{Thu05} gives exact sequences
\begin{align}
 0&\longrightarrow\mathbf R\longrightarrow\mathcal H_v
 \xrightarrow{\operatorname{div}_v}
 \operatorname{Div}(C_v)\otimes_{\mathbf Z}\mathbf R
 \xrightarrow{\deg}\mathbf R\longrightarrow0,
\label{eq:p063-harmonic-residual-exact}\\
 0&\longrightarrow\mathbf R\longrightarrow\mathcal F_v
 \xrightarrow{\operatorname{div}_v}
 \operatorname{Prin}(C_v)\otimes_{\mathbf Z}\mathbf R
 \longrightarrow0.
\label{eq:p063-smooth-residual-exact}
\end{align}
The coefficient at \(p_\xi\) is the outgoing slope \(d_\xi h(v)\).
By \eqref{eq:p063-residual-degree-zero}, the divisor \(D\) lies in the
image of the first residual map.  Its class in
\(\operatorname{Pic}^0(C_v)\otimes_{\mathbf Z}\mathbf R\) vanishes
because it is a rational multiple of a principal divisor.  Comparing
\eqref{eq:p063-harmonic-residual-exact} and
\eqref{eq:p063-smooth-residual-exact} gives a germ in \(\mathcal F_v\)
with residual divisor \(D\); adjusting by a constant gives a germ \(h\)
with \(h(v)=c\).
The construction in the proof of Thuillier's lemma permits a
representative containing short arms in every one of the finitely many
displayed directions.

By \cite[Remark~15.10]{GJR26}, \(\mathcal F\) is exactly the sheaf of
ordinary Chambert--Loir--Ducros smooth harmonic functions on a smooth
boundaryless analytic curve.  Hence \(h\) has all the properties in
\eqref{eq:p063-prescribed-residual-slopes}.
\end{proof}

\begin{proof}[Proof of
\cref{prop:p063-tree-lisse-ordinary-regularization}]
By \cite[Proposition~3.2.4]{Thu05}, after a semistable refinement and a
finite metric subdivision, there are a finite metric tree \(\Gamma\), a
retraction \(\tau:Y\to\Gamma\), and a continuous edgewise affine function
\(F\) such that
\[
                         f=F\circ\tau.
\]
Here no extension of scalars is needed because \(K\) is algebraically
closed.  Include every breakpoint and branch point among the metric
vertices, and let \(\partial\Gamma\) contain the vertices over the Shilov
boundary of \(Y\).  Interior vertices are of type \(2\) or type \(3\).
Subharmonicity and \cite[Proposition~3.4.4]{Thu05} give
\[
                         \Delta_\Gamma F(v)\geq0
                         \qquad
                         \bigl(v\notin\partial\Gamma\bigr).
\]
If \(\Gamma\) has no edge, then \(f\) is constant and there is nothing to
prove.  Otherwise apply
\cref{lem:p063-tree-rational-strictification}, with error
\(\varepsilon/3\), to obtain a rational-slope function \(G\) satisfying
\begin{equation}
\label{eq:p063-tree-graph-error}
 \lVert G\circ\tau-f\rVert_{L^\infty(Y)}
 <\frac{\varepsilon}{3}
\end{equation}
and
\begin{equation}
\label{eq:p063-tree-positive-vertex-mass}
                         \Delta_\Gamma G(v)>0
                         \qquad
                         \bigl(v\notin\partial\Gamma\bigr).
\end{equation}

We first modify \(G\circ\tau\) near each interior type-\(2\) vertex
\(v\).  Choose pairwise disjoint simple star neighbourhoods
\(U_v\Subset Y^\circ\).  Such a neighbourhood contains short arms in all
graph directions, and all but finitely many other tangent directions are
open-disc components contained in \(U_v\).  Let \(E_v\) be the finite set
of exceptional directions.  Enlarge
\[
 M_v=\{p_{v,e}:e\ni v\}\subset C_v(\widetilde K)
\]
by the residual points corresponding to \(E_v\).  This is the finite
forbidden set; in particular it contains every graph direction and every
direction in which a fresh collar is not available.

Define
\[
                         B_v
   =\sum_{e\ni v}d_eG(v)[p_{v,e}]
   \in\operatorname{Div}_{\mathbf Q}(C_v).
\]
By \eqref{eq:p063-tree-positive-vertex-mass},
\[
                         \deg B_v
                         =\Delta_\Gamma G(v)>0.
\]
Apply \cref{lem:p063-residual-divisor-moving} to \(B_v\) and \(M_v\).
There are \(N_v>0\) and an effective divisor
\[
                         A_v=\sum_{p\in S_v}m_{v,p}[p]
\]
such that
\[
 A_v\sim N_vB_v,
 \qquad
 S_v\cap M_v=\varnothing.
\]
Thus
\[
                         D_v
    =B_v-\frac1{N_v}A_v
\]
has degree zero and is a rational multiple of a principal divisor.  By
\cref{lem:p063-principal-residual-germ}, there is an ordinary smooth
harmonic germ \(h_v\), normalized by \(h_v(v)=G(v)\), whose graph slopes
are \(d_eG(v)\), whose slopes in the fresh directions \(p\in S_v\) are
\[
                         -a_{v,p},
             \qquad
                         a_{v,p}=\frac{m_{v,p}}{N_v}>0,
\]
and whose slopes in all remaining directions vanish.

We now choose one representative on one neighbourhood, rather than making
direction-by-direction choices.  Write it as a finite real linear
combination of logarithms of analytic units.  After one common shrink, each
unit has a rational reduction on \(C_v\); the union \(T_v\) of the supports
of their divisors is finite.  On every reduction fibre outside \(T_v\),
the absolute value of every unit is constant.  In a direction belonging to
\(T_v\) but not to the support of \(D_v\), the sum has outgoing slope zero;
after shrinking the finitely many such annular collars, its affine germ is
therefore constant.  This uses the reduction-fibre description of a simple
type-\(2\) neighbourhood in
\cite[\S4.29 and Lemma~5.12(3)]{BPR13Structure}.

Shrink once more along the finitely many graph and fresh directions.  On
this single neighbourhood, \(h_v\) agrees with \(G\circ\tau\)
simultaneously on every graph collar and on every zero-slope overlap: on a
graph collar the two functions have the same value and affine slope, and on
all remaining directions both are constant with value \(G(v)\).  The
avoidance of \(E_v\) ensures that each \(p\in S_v\)
determines an available open-disc component of \(Y\setminus\{v\}\),
disjoint from \(\Gamma\).

We cap the negative slope in each of these fresh directions.  Fix an even
nonnegative function \(\rho\in C^\infty_c((-1,1))\) of integral one and
put
\[
 \chi_\eta(r)=\int_{\mathbf R}|r-\eta s|\rho(s)\,\mathrm ds,
 \qquad
 m_\eta(s,t)=\frac{s+t+\chi_\eta(s-t)}2.
\]
Then \(m_\eta\) is smooth, convex, and nondecreasing in each variable,
\begin{equation}
\label{eq:p063-regularized-max-bounds}
 \max\{s,t\}\leq m_\eta(s,t)
       \leq\max\{s,t\}+\frac{\eta}{2},
\end{equation}
and it equals \(\max\{s,t\}\) if \(|s-t|\geq\eta\).  Indeed, Jensen's
inequality gives \(\chi_\eta(r)\geq|r|\), the triangle inequality gives
\(\chi_\eta(r)\leq|r|+\eta\), evenness gives zero first moment, and hence
\(\chi_\eta(r)=|r|\) for \(|r|\geq\eta\).  Also
\(|\chi_\eta'|\leq1\) and \(\chi_\eta''\geq0\).

Fix \(p\in S_v\).  On a sufficiently short punctured neighbourhood in its
open-disc component, choose an invertible annular coordinate \(T\) and a
constant \(c\) such that
\[
                         t=-\log|T|+c\geq0
\]
is distance from \(v\) on the skeleton.  Shrinking the whole annular tube,
not merely its skeletal arm, we have
\[
                         h_v=G(v)-a_{v,p}t.
\]
Choose \(r_{v,p}>0\) and \(\eta_{v,p}>0\) so that
\[
 \eta_{v,p}<a_{v,p}r_{v,p},
 \qquad
 r_{v,p}+\frac{\eta_{v,p}}{a_{v,p}}
\]
lies inside the arm on which this formula holds, and
\begin{equation}
\label{eq:p063-fresh-collar-error-choice}
 a_{v,p}r_{v,p}+\frac{\eta_{v,p}}2
 <\frac{\varepsilon}{3}.
\end{equation}
On that entire annular tube set
\begin{equation}
\label{eq:p063-fresh-collar-cap}
 q_{v,p}(t)
  =
 m_{\eta_{v,p}}\bigl(
     G(v)-a_{v,p}t,\,
     G(v)-a_{v,p}r_{v,p}
   \bigr).
\end{equation}
It equals \(h_v\) on
\[
               0\leq t<
               r_{v,p}-\frac{\eta_{v,p}}{a_{v,p}}
\]
and equals the constant \(G(v)-a_{v,p}r_{v,p}\) when
\[
               t>
               r_{v,p}+\frac{\eta_{v,p}}{a_{v,p}}.
\]
The same formula therefore assigns values on all residue discs attached to
the skeletal arm.  Extend it by that constant over the rest of the
open-disc component in the direction \(p\).  No return to the old constant
\(G(v)\) is made or needed.  By
\eqref{eq:p063-regularized-max-bounds},
\begin{equation}
\label{eq:p063-fresh-collar-error}
 \left|q_{v,p}-G(v)\right|
 \leq a_{v,p}r_{v,p}+\frac{\eta_{v,p}}2
 <\frac{\varepsilon}{3}.
\end{equation}

The function in \eqref{eq:p063-fresh-collar-cap} is ordinary smooth and
psh.  Indeed, \(t\) differs by a constant from \(-\log|T|\) for an
invertible analytic annular coordinate, and
\eqref{eq:p063-fresh-collar-cap} is a smooth convex function of \(t\) on
the whole annular tube.
Equivalently, it is a smooth convex coordinatewise nondecreasing
regularized maximum of the ordinary smooth harmonic functions \(h_v\)
and \(G(v)-a_{v,p}r_{v,p}\).  Near \(v\) every cap equals \(h_v\), while
at its far end it equals its constant extension on an open overlap.
Since \(S_v\) is finite, the modifications therefore glue around \(v\)
to an ordinary smooth psh function.

It remains to smooth the corners at interior type-\(3\) subdivision
vertices.  A smooth analytic curve has exactly two tangent directions at
such a point.  In an annular coordinate \(t\), with the vertex at \(0\),
write
\[
 G(t)=
 \begin{cases}
  G(v)+\alpha t,&t\leq0,\\
  G(v)+\beta t,&t\geq0.
 \end{cases}
\]
The outgoing-slope sum is \(\beta-\alpha\), so
\eqref{eq:p063-tree-positive-vertex-mass} gives \(\alpha<\beta\), and
\[
 G(t)=\max\{G(v)+\alpha t,G(v)+\beta t\}.
\]
Choose a collar on which the two affine formulas hold and which is disjoint
from every other replacement collar.  On that collar replace this maximum
by
\[
 m_\eta\bigl(G(v)+\alpha t,G(v)+\beta t\bigr).
\]
Choose \(\eta>0\) so small that
\(\eta/(\beta-\alpha)\) lies strictly inside this chosen collar.  The
replacement agrees with \(G\) whenever
\[
                         |t|>\frac{\eta}{\beta-\alpha},
\]
is ordinary smooth and psh, and differs from \(G\) by at most
\(\eta/2\).  Decrease \(\eta\), if necessary, so that also
\(\eta<2\varepsilon/3\).

There are only finitely many interior type-\(2\) and type-\(3\) vertices,
so all graph collars can be chosen pairwise disjoint.  At a type-\(2\)
vertex the central function is \(h_v\); on every graph or zero-slope arm
it equals \(G\circ\tau\) on an open overlap.  Every fresh direction is
off the graph, different fresh directions give disjoint open-disc
components, and the cap equals its constant continuation on an open
overlap.  The type-\(3\) replacement equals \(G\circ\tau\) on open
overlaps at both ends.  Away from these collars,
\(G\circ\tau\) is either affine in an ordinary annular coordinate or
locally constant.  Thus all local functions glue to a continuous
function \(g\) on \(Y\), and \(g|_{Y^\circ}\) is ordinary smooth and psh.
Boundary vertices need no smoothing because smoothness is asserted only
on \(Y^\circ\).

The local replacement supports are disjoint.  By
\eqref{eq:p063-fresh-collar-error} and the type-\(3\) error bound,
\[
                \lVert g-G\circ\tau\rVert_{L^\infty(Y)}
                    <\frac{\varepsilon}{3}.
\]
Together with \eqref{eq:p063-tree-graph-error}, this gives
\[
                         \lVert g-f\rVert_{L^\infty(Y)}
                         <\frac{2\varepsilon}{3}<\varepsilon,
\]
which proves \eqref{eq:p063-lisse-uniform-regularization}.
\end{proof}

\begin{proof}[Proof of
\cref{thm:p063-current-psh-ordinary-regularization}]
The equivalence between
\eqref{eq:p063-current-positivity} and Thuillier subharmonicity for
continuous functions is \cite[Corollary~4.8]{Wanner19}.
We prove that condition \(2\) implies condition \(3\).

Fix \(x\in W\).  By \cref{lem:p063-local-tree-affinoid}, choose a connected
strict affinoid tree domain \(Y\subset W\) with \(x\in Y^\circ\).  Choose
an open set \(\Omega\) with
\[
                         Y\subset\Omega\Subset W.
\]
Thuillier's approximation theorem
\cite[Theorem~3.4.2]{Thu05} gives a decreasing net
\((f_i)_i\) of lisse subharmonic functions on \(\Omega\), converging
pointwise to \(u\).  Its restriction to the compact space \(Y\) converges
uniformly.  Indeed, for \(\delta>0\), the increasing family of open sets
\[
                         U_i=\{y\in Y:f_i(y)-u(y)<\delta\}
\]
covers \(Y\).  A finite subcover and directedness provide one index \(i\)
with \(U_i=Y\).  This is the net form of Dini's argument; it uses the
continuity of \(u\) and of every \(f_i\).

Put \(b_n=4^{-n}\).  Recursively choose an increasing sequence of indices
\(i_n\) such that
\[
                         \lVert f_{i_n}-u\rVert_{L^\infty(Y)}
                         <\frac{b_n}{2}.
\]
At the \(n\)-th step, take a common upper bound of \(i_{n-1}\) and an
index supplied by uniform convergence.  Apply
\cref{prop:p063-tree-lisse-ordinary-regularization} to
\(f_{i_n}|_Y\), with error \(b_n/2\).  We obtain \(g_n\in C(Y)\), ordinary
smooth and psh on \(Y^\circ\), with
\begin{equation}
\label{eq:p063-raw-uniform-error}
                         \lVert g_n-u\rVert_{L^\infty(Y)}<b_n.
\end{equation}

The functions \(g_n\) need not be ordered.  Define
\begin{equation}
\label{eq:p063-tail-shift}
                         c_n
   =\sum_{j=n}^{\infty}(b_j+b_{j+1}),
   \qquad
                         u_n=g_n+c_n.
\end{equation}
Since \(c_n-c_{n+1}=b_n+b_{n+1}\),
\eqref{eq:p063-raw-uniform-error} gives, pointwise on \(Y\),
\[
 \begin{aligned}
 u_{n+1}-u_n
   &=
 g_{n+1}-g_n-(b_n+b_{n+1})\\
   &\leq
 |g_{n+1}-u|+|g_n-u|-(b_n+b_{n+1})
 \leq0.
 \end{aligned}
\]
Thus \(u_{n+1}\leq u_n\).  Moreover,
\[
                         \lVert u_n-u\rVert_\infty
                         \leq b_n+c_n\longrightarrow0.
\]
Adding constants preserves ordinary smoothness and psh.  This proves
\eqref{eq:p063-monotone-ordinary-approximation}.

Finally assume condition \(3\).  On \(Y^\circ\), each
\(\mathrm d'\mathrm d''[u_n]\) is a positive current.  Uniform convergence
gives \([u_n]\to[u]\) as currents: for every compactly supported smooth
test function \(\varphi\), the absolute value of
\[
 \bigl\langle [u_n-u],\mathrm d'\mathrm d''\varphi\bigr\rangle
\]
is at most \(\lVert u_n-u\rVert_\infty\) times the total variation of the
compactly supported smooth signed measure
\(\mathrm d'\mathrm d''\varphi\).  Hence
\[
                    \mathrm d'\mathrm d''[u_n]
                    \longrightarrow
                    \mathrm d'\mathrm d''[u]
\]
weakly.  Positivity is closed under weak convergence, so
\(\mathrm d'\mathrm d''[u]\geq0\) on \(Y^\circ\).  These interiors cover
\(W\), proving condition \(1\).

Condition \(3\) is exactly local membership, for continuous functions, in
the first decreasing-limit stage generated by ordinary smooth psh
functions.  The convergence obtained above is even uniform.
\end{proof}
